\documentclass[11pt,a4paper]{article}

\usepackage[T1]{fontenc}
\usepackage{amsmath,amssymb,amsthm,mathtools}
\usepackage{booktabs}
\usepackage{enumitem}
\usepackage{geometry}
\usepackage{microtype}
\usepackage{listings}
\usepackage{hyperref}

\geometry{margin=25mm}
\hypersetup{
  colorlinks=true,
  linkcolor=blue,
  citecolor=blue,
  urlcolor=blue,
  pdftitle={Maximum Determinants of 55 by 55 Circulant Matrices over 0,1 and -1,1},
  pdfauthor={Qichao Wang}
}
\urlstyle{same}
\setlength{\parindent}{0pt}
\setlength{\parskip}{5pt}
\setlength{\emergencystretch}{2em}
\lstset{
  basicstyle=\ttfamily\small,
  columns=fullflexible,
  keepspaces=true,
  breaklines=true,
  frame=single,
  framerule=0.3pt,
  xleftmargin=4pt,
  xrightmargin=4pt
}

\newtheorem{theorem}{Theorem}
\newtheorem{proposition}{Proposition}
\newtheorem{lemma}{Lemma}
\theoremstyle{remark}
\newtheorem*{remark}{Remark}
\DeclareMathOperator{\Res}{Res}
\DeclareMathOperator{\wt}{wt}
\newcommand{\Z}{\mathbb Z}
\newcommand{\F}{\mathbb F}
\newcommand{\D}{D_{01}(55)}
\newcommand{\M}{134694094094758395331307111329132}

\title{Maximum Determinants of $55\times55$ Circulant Matrices\\over $\{0,1\}$ and $\{-1,1\}$}
\author{
  Qichao Wang\\
  Hebei University of Technology
}
\date{}

\begin{document}
\maketitle

\begin{abstract}
We determine the exact maximum determinants of order-$55$ circulant
matrices over $\{0,1\}$ and $\{-1,1\}$ and classify all maximizing
words.  The binary maximum is
$134694094094758395331307111329132$; its maximizers have weight $28$
and form one affine orbit of size $2200$.  The sign maximum is
$297532404423965431213849494795385191907933028352$; its
maximizers form two affine orbits of size $2200$, exchanged by global
negation.  The common Fourier spectrum yields two distinct determinant
objectives through their zero-frequency factors.  Exact moment
bounds, cyclotomic folds, correlation profiles, and exhaustive binary
lifts establish both extrema and their complete classifications.
\end{abstract}

\section{Introduction}

For $a\in\{0,1\}^{55}$, let $C(a)$ denote its circulant
matrix.  The two problems considered here are
\[
 D_{01}(55)=\max_{a\in\{0,1\}^{55}}|\det C(a)|,\qquad
 D_{\pm}^{\mathrm{circ}}(55)=
       \max_{x\in\{-1,1\}^{55}}|\det C(x)|.
\]
Exact binary-circulant tables of Brent and Yedidia extend through
order $53$; an order-$54$ computation supplies the next reference
point \cite{brent-yedidia,aletheia54}.
At $55=5\cdot11$, the discrete Fourier spectrum separates into
cyclotomic sectors of orders $1,5,11,55$, enabling a common
finite reduction for the two alphabets.

\begin{theorem}[Global maximum and classification]
\label{thm:main}
For binary circulant matrices of order $55$,
\[
  D_{01}(55)=M=\M.
\]
Define $w=(w_0,\ldots,w_{54})$ by its digits, read from left to right:
\begin{center}
\texttt{0000000100011111011011001101011011010100011110101000111}
\end{center}
It has weight $28$ and attains $M$. Every maximizing word is of the form
\[
 (w_{ui+t})_{i=0}^{54},
 \qquad u\in(\Z/55\Z)^\times,\quad t\in\Z/55\Z.
\]
The $2200$ words in this affine orbit are distinct.
\end{theorem}

\begin{theorem}[Sign maximum and classification]\label{thm:sign}
For sign circulants of order $55$,
\[
 D_{\pm}^{\mathrm{circ}}(55)=
 297532404423965431213849494795385191907933028352,
\]
\[
 S_{\pm}(55)=D_{\pm}^{\mathrm{circ}}(55)/2^{54}
 =16516366298178510386365752382353.
\]
Define $v=(v_0,\ldots,v_{54})\in\{0,1\}^{55}$ by
\begin{center}
\texttt{0000000000100110110100111001001110001010101100010010111}
\end{center}
Every maximizing sign word has the form
\[
  x_i=\varepsilon(2v_{ui+t}-1),\qquad
  \varepsilon\in\{-1,1\},\quad
  u\in(\Z/55\Z)^\times,\quad
  t\in\Z/55\Z.
\]
These are $4400$ distinct words in two affine orbits of $2200$
each.  Their plus-sign counts and row sums are $(23,-9)$ and
$(32,9)$, respectively.
\end{theorem}

The two determinant objectives share the product of the $27$
nontrivial Fourier magnitudes.  Their zero-frequency coefficients
differ: after complement reduction the binary objective has factor
$55-k$, and the normalized sign objective has factor $55-2k$,
where $k$ is the number of ones in a binary word.  Thus the
same spectral bounds lead to distinct active weights and distinct
maximizing words.

The proof develops Fourier and autocorrelation identities, an
exact binary-to-sign relation, moment and cyclotomic bounds, and
complete enumeration of compatible binary torus lifts.
The two classifications follow from separate exact finite
certificates, with their arithmetic verification collected in the
appendix.

\section{Fourier formulation, symmetry and the sign bridge}

\subsection{Binary Fourier spectrum and autocorrelation}


Let $a\in\{0,1\}^{55}$, let $k=\sum_{j=0}^{54}a_j$, and define
\[
  C(a)_{ij}=a_{j-i},
  \qquad
  f_a(x)=\sum_{j=0}^{54}a_jx^j.
\]
For $\zeta=\exp(2\pi i/55)$, the Fourier eigenvalues of
$C(a)$ are $f_a(\zeta^r)$, $0\leq r<55$.  Since $a$ is real and $55$ is odd, the nonzero frequencies
pair by complex conjugation.  Thus, with
\[
  q_r=|f_a(\zeta^r)|^2,\qquad 1\leq r\leq27,
\]
we have
\begin{equation}
\label{eq:fourier-det}
  \det C(a)=k\prod_{r=1}^{27}q_r\geq0.
\end{equation}
Every binary determinant is consequently nonnegative.  Formal
correlation profiles are screened using absolute spectral products
before binary realizability is established.

Define the periodic autocorrelation sequence
\[
  c_s=\sum_{t=0}^{54}a_ta_{t+s},
  \qquad s\in\Z/55\Z.
\]
Then $c_0=k$, $c_s=c_{55-s}$, and counting ordered pairs of distinct
support elements gives
\begin{equation}
\label{eq:corr-sum}
  2\sum_{s=1}^{27}c_s=k(k-1).
\end{equation}
The correlation sequence is the inverse Fourier transform of the squared
Fourier magnitudes.  Parseval therefore gives the two moment identities
\begin{align}
  2\sum_{r=1}^{27}q_r&=55k-k^2, \label{eq:first-moment}\\
  k^4+2\sum_{r=1}^{27}q_r^2
    &=55\left(k^2+2\sum_{s=1}^{27}c_s^2\right).
    \label{eq:fourth-moment}
\end{align}
The first identity and AM--GM immediately imply
\begin{equation}
\label{eq:amgm}
  \det C(a)\leq k\left(\frac{k(55-k)}{54}\right)^{27}.
\end{equation}


\subsection{Complementation and affine symmetry}


Let $\bar a=1-a$.  At the zero frequency,
$f_{\bar a}(1)=55-k$.  At every nonzero $55$th root of unity,
\[
  f_{\bar a}(\zeta^r)=-f_a(\zeta^r),
  \qquad 1\leq r<55,
\]
because the sum of all $55$th roots in a nontrivial frequency is zero.
There are $54$ nonzero eigenvalues, so the signs cancel in the determinant.
For $0<k<55$ this proves
\begin{equation}
\label{eq:complement}
  k\det C(\bar a)=(55-k)\det C(a).
\end{equation}
The constant words are singular.  If
\[
  D_{\mathrm{circ}}(55,k)=
  \max_{\wt(a)=k}\det C(a),
\]
then complement is a bijection between weight strata and
\begin{equation}
\label{eq:weight-pair}
  D_{\mathrm{circ}}(55,55-k)
    =\frac{55-k}{k}D_{\mathrm{circ}}(55,k).
\end{equation}
It is therefore enough to handle $1\leq k\leq27$.  For the complementary
word, the first-moment bound takes the form
\begin{equation}
\label{eq:complement-amgm}
  \det C(\bar a)
    =(55-k)\prod_{r=1}^{27}q_r
    \leq(55-k)\left(\frac{k(55-k)}{54}\right)^{27}.
\end{equation}

Translations and multipliers preserve the determinant.  Translation changes
Fourier eigenvalues only by phases whose product is one, while multiplication
of indices by $u\in(\Z/55\Z)^\times$ permutes the nonzero frequencies.  The
affine group has order
\[
  55\varphi(55)=55\cdot40=2200.
\]
Reversal is the multiplier $u=-1$.  Orbit sizes follow from the stabilizers of their words.
Correlations are translation-invariant and are indexed
by the $27$ opposite-shift pairs.  On these pairs the effective unit action
has order $20$.


\subsection{Binary-to-sign determinant identity}

\begin{lemma}[Binary-to-sign bridge]\label{lem:sign-bridge}
For $a\in\{0,1\}^{55}$ of weight $0<k<55$ and $x=2a-\mathbf1$,
\[
 \det C(x)=2^{54}\frac{2k-55}{k}\det C(a),\qquad
 S_{\pm}(55)=\max_{1\le k\le27}\frac{55-2k}{k}D_{\mathrm{circ}}(55,k).
\]
\end{lemma}
\begin{proof}
The identity $C(x)=2C(a)-J$ changes the eigenvalue on the all-ones
line to $2k-55$ and doubles every nontrivial Fourier eigenvalue.
It also holds when the circulant is singular. Conjugate pairing gives
$\det C(a)=k\prod q_r$. Complementing $a$ negates $x$ and preserves
absolute determinant. The constant endpoints have rank one.
\end{proof}
Thus the reduced objective is
\[
 R(a)=(55-2k)\prod_{r=1}^{27}q_r=(55-2k)N_5N_{11}N_{55}.
\]
For the encoding $x=2a-\mathbf1$, the signed row sum is $2k-55$;
$55-2k$ is its magnitude on the reduced weights.
In contrast, the complementary binary objective has coefficient $55-k$.
The small-weight results in Appendix~\ref{app:small}
transfer through the bridge. At reduced weight $27$, binary complement duality gives
the maximum normalized sign value
\[
 M/28=4810503360527085547546682547469.
\]



\section{Moment and cyclotomic reduction}
\label{sec:reduction}

\subsection{Moment bounds and exact spectral products}

The correlation identities turn the determinant problem into an
integer moment problem.  Among $h$ nonnegative integers with fixed sum
$T$, write $T=hb+r$ with $0\leq r<h$.  Their minimum square sum is
\[
  (h-r)b^2+r(b+1)^2.
\]
Transferring one unit between entries that differ by at least two
decreases their square sum.  Applied to the $27$ independent
correlations, this gives the optimal lower bound on
$\sum_{s=1}^{27}c_s^2$ at each weight.

The fourth moment in \eqref{eq:fourth-moment} fixes the square sum of the
$q_r$.  The following elementary extremal calculation supplies a rigorous
upper bound on their product.  Consider $h$ nonnegative variables with sum
$S$ and square sum at least $L$.  Contracting variables toward their common
mean increases the product until the square-sum constraint is tight (or the
variance is zero).  At a positive constrained stationary point, Lagrange
multipliers give
\[
  \frac1x=\alpha+2\beta x,
\]
so there are at most two distinct values.  If the smaller value occurs $m$
times, put $\mu=S/h$ and $V=L-S^2/h$.  The two values are
\begin{equation}
\label{eq:two-values}
  \mu-\sqrt{\frac{V(h-m)}{hm}},
  \qquad
  \mu+\sqrt{\frac{Vm}{h(h-m)}}.
\end{equation}
Enumerating $m=1,\ldots,h-1$, together with the zero-variance and boundary
cases, exhausts the compact feasible set.  The implementation bounds the
square roots outward by integer arithmetic at a decimal scale of $10^{50}$.
Every pruning decision uses exact integers and outward rational bounds.

The certificate also imposes a uniform spectral cap.  Let
\[
  b=\left\lfloor\frac{k(k-1)}{54}\right\rfloor,
  \qquad \delta_s=c_s-b.
\]
For every nonzero frequency,
\begin{equation}
\label{eq:spectral-cap-formula}
  q_r=k-b+\sum_{s=1}^{27}\delta_s\,2\cos(2\pi rs/55).
\end{equation}
The rearrangement inequality pairs the sorted deviations with sorted cosine
coefficients.  There are only three coefficient multisets, represented by
$r=1,5,11$, according to the gcd class of the frequency.  Their rational
outward bounds give a cap $B$ valid for every $q_r$.

To maximize a product subject also to $0\leq x_i\leq B$, enumerate the
number of coordinates on the upper boundary.  The remaining interior
coordinates again have the form \eqref{eq:two-values}; infeasible sums,
negative roots, and roots above $B$ are discarded.  This is a finite
description of all extrema on the capped set.  The cosine intervals are
constructed from Machin's identity and rational Taylor remainders for
$\pi$ and cosine, and the stored dyadic intervals are independently checked
by the verifier.


\subsection{Cyclotomic folding modulo $5$ and $11$}

The factorization $55=5\cdot11$ splits the frequency space into the
cyclotomic sectors of orders $1,5,11,55$.  Their dimensions are
\[
  1,\quad4,\quad10,\quad40,
  \qquad 1+4+10+40=55.
\]
The nonzero sectors have respectively $2$, $5$, and $20$ conjugate pairs.
Let
\[
  N_m(a)=\Res(f_a,\Phi_m),
  \qquad m\in\{5,11,55\}.
\]
All three cyclotomic degrees are even, so the resultant convention used here
has positive sign, and conjugate pairing makes these norms nonnegative.  We
obtain
\begin{equation}
\label{eq:resultant-factorization}
  \det C(a)=kN_5(a)N_{11}(a)N_{55}(a).
\end{equation}

Use the Chinese remainder identification
\[
  \Z/55\Z\cong\Z/5\Z\times\Z/11\Z.
\]
For $m\in\{5,11\}$ define the folded vector
\[
 F_i^{(m)}=\sum_{t\equiv i\pmod m}a_t,\qquad 0\leq i<m.
\]
Write $r=F^{(5)}$ and $s=F^{(11)}$, so
$0\leq r_i\leq11$ and $0\leq s_j\leq5$.
The norms $N_5$ and $N_{11}$ depend only on their respective folds.
The folded autocorrelations satisfy
\[
 H_i^{(m)}=\sum_{j=0}^{m-1}F_j^{(m)}F_{j+i}^{(m)}
          =\sum_{t\equiv i\pmod m}c_t.
\]
For $\zeta_m=\exp(2\pi i/m)$, put
$Q_j^{(m)}=|f_a(\zeta_m^j)|^2$.  The first two power sums
over the $(m-1)/2$ nonzero conjugate pairs are
\begin{align}
 S_1^{(m)}&=\sum_{j=1}^{(m-1)/2}Q_j^{(m)}
   =\frac{m\sum_{i=0}^{m-1}(F_i^{(m)})^2-k^2}{2},\\
 S_2^{(m)}&=\sum_{j=1}^{(m-1)/2}(Q_j^{(m)})^2
   =\frac{m\sum_{i=0}^{m-1}(H_i^{(m)})^2-k^4}{2}.
\end{align}

For fixed folds $r$ and $s$, the first moment in the primitive order-$55$
sector is
\begin{equation}
\label{eq:primitive-sector-sum}
  T=\frac{55k+k^2-5\sum_i r_i^2-11\sum_j s_j^2}{2}.
\end{equation}
Its norm is at most $(T/20)^{20}$.  A simpler necessary fold screen is also
useful.  If the selected small sector has $h=(m-1)/2$ pairs and sum $T_m$,
the remaining $27-h$ pairs have sum
\[
  R=\frac{k(55-k)}2-T_m.
\]
For the binary objective, a complementary determinant of at least $M$
requires
\begin{equation}
\label{eq:fold-screen}
  (55-k)N_m\left(\frac{R}{27-h}\right)^{27-h}\geq M.
\end{equation}
This is rational and exhaustive over the bounded integer fold vectors.

All unit orientations of the folded correlations are retained.  Independent
translations of the two fold vectors lift jointly by the Chinese remainder
theorem.  The binary lifts of all surviving fold pairs are enumerated in
Section~\ref{sec:lifts}.

\subsection{The common finite decision chain}

For either determinant objective, the exact upper bounds reduce
the correlation and fold data to a finite list of targets.
Each target becomes a binary matrix enumeration on
$\mathbb Z/5\mathbb Z\times\mathbb Z/11\mathbb Z$.
The two objective coefficients are
\[
 \begin{array}{ll}
 (55-k)\displaystyle\prod_{r=1}^{27}q_r,
   &\hbox{binary complement objective},\\[4pt]
 (55-2k)\displaystyle\prod_{r=1}^{27}q_r,
   &\hbox{normalized sign objective}.
 \end{array}
\]
The same spectrum bounds both objectives, with thresholds set by
their respective incumbents.  The following sections give the
complete target lists, lift counts and equality cases.

\section{The binary maximum}\label{sec:binary}

\subsection{Active weights and folded profiles}

The moment and exact product bounds exclude every lower
complementary weight $k\le24$ at the incumbent $\M$.
The surviving correlation-energy windows are shown below.
Each interval bounds $\sum_{s=1}^{27}c_s^2$; the upper endpoint is
the largest integer compatible with reaching $M$ under the
certified moment bound.
\begin{center}
\small
\begin{tabular}{@{}rrrr@{}}
\toprule
lower weight $k$ & complement weight & minimum square sum & maximum square sum\\
\midrule
25 & 30 & 3336 & 3338\\
26 & 29 & 3913 & 3919\\
27 & 28 & 4563 & 4570\\
\bottomrule
\end{tabular}
\end{center}
For $k=27$, the deviations $d_s=c_s-13$ have sum zero, so
$\sum d_s^2$ is even; the final row is equivalently a small deviation
window around the nearly balanced correlation $13$.


The mod-$5$ and mod-$11$ screens retain the following numbers of
representatives under their affine fold symmetries:
\begin{center}
\small
\begin{tabular}{@{}rrr@{}}
\toprule
$k$ & order-$5$ representatives & order-$11$ representatives\\
\midrule
25 & 17 & 321\\
26 & 27 & 560\\
27 & 26 & 660\\
\bottomrule
\end{tabular}
\end{center}

\subsection{Correlation profiles and spectral screening}


For each active weight $k\in\{25,26,27\}$, the exact sum in
\eqref{eq:corr-sum}, the square-sum window, the range $0\leq c_s\leq k$, and
the capped moment bound define a finite family of integer correlation
profiles.  The generator first visits sorted multisets and uses the
multinomial coefficient to recover the exact number of ordered profiles.
It then tests every coordinate assignment against both fold signatures and
the effective order-$20$ unit action.

The table records the exact certificate counts.
``Allowed'' counts formal ordered profiles from the initial integer
screen, and ``explicitly enumerated'' gives the completed exact
evaluations.  Binary lifts decide realizability of retained profiles.
\begin{center}
\scriptsize
\setlength{\tabcolsep}{4pt}
\begin{tabular}{@{}lrrrr@{}}
\toprule
$k$ & allowed ordered profiles & explicitly enumerated & canonical profiles & lift targets\\
\midrule
25 & 407,277 & 2,925 & 55 & 55\\
26 & 32,178,654 & 816,102 & 22,375 & 281\\
27 & 6,043,753 & 6,026,203 & 155,558 & 1,980\\
\midrule
total & 38,629,684 & 6,845,230 & 177,988 & 2,316\\
\bottomrule
\end{tabular}
\end{center}

For a formal profile, the $q_r$ obtained from the discrete Fourier transform
can be negative.  The profile screen therefore applies the absolute-product
inequality
\begin{equation}
\label{eq:formal-profile-bound}
  \left|(55-k)\prod_{r=1}^{27}q_r\right|
  \leq (55-k)\left(\frac{\sum_r q_r^2}{27}\right)^{27/2},
\end{equation}
which is just Cauchy--Schwarz followed by AM--GM on the absolute values.
The signed product is reconstructed from two $61$-bit Fourier fields.
A squared integer comparison establishes that the combined CRT modulus
exceeds the upper bound in \eqref{eq:formal-profile-bound}.
The reconstruction therefore preserves the exact product.
Profiles with absolute product below $M$ are eliminated; the remaining
profiles enter the binary lift enumeration.


\subsection{Exhaustive binary lifts}\label{sec:lifts}

Use the explicit torus coordinate
\[
  (i,j)\in\Z/5\Z\times\Z/11\Z
  \longmapsto 11i+45j\pmod{55}.
\]
The $11$ columns are five-bit masks, so there are $32$ possible column
states.  For every retained correlation target, all compatible oriented
pairs of fold vectors are enumerated after normalizing the independent row
and column translations.  This produces exactly $16,084$ lift jobs.

The production lift splits the eleven columns into five and six.  A partial
assignment records its word and a seven-coordinate signature: five row
counts and the two within-column correlation sums for shifts $11$ and $22$.
These statistics are additive across columns.  The left half is sorted by
its packed signature; each right half is joined to the unique complementary
signature.  The packed signature is injective over the bounded coordinate ranges,
so each join uses exact integer equality.  Every joined word has the
requested folds and stored correlations.  The remaining shifts are
checked by rotated $55$-bit intersections and population counts.

A partial assignment is discarded when any nonnegative row count or
correlation exceeds its target.  The production five-plus-six join tests
$15,487,882,832$ complete words and returns two normalized solutions.
An independently implemented six-plus-five join uses typed keys and
exhausts all $16,084$ tasks.  Both implementations agree on every
task's join count and output words.
The complete supplement is
\path|certificates/order55_sign/binary_companion/full_fiber_report.json|.
A separate small fiber is checked against all $78,125$ direct column
assignments using three splits.


\subsection{Coverage theorem and exact enumeration}

\begin{lemma}[Complete binary search]\label{lem:binary-complete}
For every binary word $a$ with $\det C(a)\geq M$, either $a$ or
its complement has an affine representative in one of the $16,084$
reduced-weight lift tasks.  Both exhaustive implementations return
the same two normalized solutions.
\end{lemma}
\begin{proof}
Complementation reduces the candidate weights to $k\le27$.
The moment and product bounds exclude $k\le24$.
For $k=25,26,27$, the correlation-sum and square-sum conditions
produce the finite integer multiset families.  The folded-product
screens, unit action and exact profile comparisons cover every
eligible ordered correlation by one of the $2,316$ target records.

For each target, every compatible mod-$5$ and mod-$11$ margin
orientation is enumerated.  Chinese remaindering realizes
independent fold rotations as one translation modulo $55$.
The split-column enumerator retains exactly the prescribed row
counts and within-column correlations; the completed word is
checked at every remaining shift.  Thus every binary realization
in the normalized target fiber appears in the enumeration.
The complete task set is exhausted by both independent implementations.
Their join counts and outputs agree for every task.
\end{proof}

The independent profile generator reconstructs the $16$ multiset
families and all $2,316$ target triples.  Its certificate entries
contain the correlation data, exact upper bounds and pruning decisions.
All threshold comparisons use integer arithmetic or outward rational bounds.



\subsection{Extremizers and their affine classification}


For the word $w$ in Theorem~\ref{thm:main}, the exact folded
data and cyclotomic norms are
\begin{lstlisting}
weight = 28
mod-5 margin:  [5,4,8,6,5]
mod-11 margin: [3,1,4,2,3,1,1,3,3,3,4]
N5  = 131
N11 = 434039
N55 = 84603917386870862636041
\end{lstlisting}
Equation \eqref{eq:resultant-factorization} gives
$28\cdot131\cdot434039\cdot84603917386870862636041=M$.
Its complement is
\begin{center}
\texttt{1111111011100000100100110010100100101011100001010111000}
\end{center}
and has weight $27$. The three nontrivial norms are unchanged, so
\[
 \det C(\bar w)=\frac{27}{28}M
  =129883590734231309783760428781663.
\]

The nonzero correlations of the winner have the distribution
\[
  13\text{ four times},\qquad
  14\text{ forty-six times},\qquad
  15\text{ four times}.
\]
Among the independent shifts $1\leq s\leq27$, the exceptional shifts are
$2,6$ with value $15$ and $7,15$ with value $13$.  For the complement, the
half-profile deviations from $13$ have square sum $4$.  If the word is
converted to signs by $x_t=2a_t-1$, its nonzero periodic autocorrelations
are $-5,-1,3$ with total squared energy $182$ over all $54$ nonzero shifts.

The perfectly balanced weight-$27$ possibility would have
$c_s=13$ for every nonzero shift, corresponding to a cyclic
$(55,27,13)$ difference set.  It is impossible.  Indeed, with
$\alpha=f_a(\zeta_5)$, where $\zeta_5$ is a primitive fifth root, the ideal
correlations imply $\alpha\overline\alpha=14$.  The polynomial
$\Phi_5$ is irreducible modulo $2$, since $\operatorname{ord}_5(2)=4$;
hence $\mathbb Z[\zeta_5]/(2)$ is a field.  Reducing the norm equation
modulo $2$ forces $\alpha=2\beta$.  Then
$\beta\overline\beta=7/2$, impossible because a rational algebraic integer
is an integer.  This obstruction also appears in the design literature
\cite{braic}.

The two normalized lift solutions have the same affine canonical
representative $\bar w$, each with complementary determinant $M$.
Lemma~\ref{lem:binary-complete} supplies the global upper bound,
and the exact norm product above shows that $w$ attains it.
Enumerating the affine images of $w$ gives $2200$ distinct words,
so its stabilizer is trivial.  These are all equality cases,
proving Theorem~\ref{thm:main}.


\section{The sign maximum}\label{sec:sign}

The exact sign optimum is recorded as OEIS A215723, and the
normalized value as A215897 \cite{oeis-sign,oeis-sign-scaled}.
Its objective is obtained from Lemma~\ref{lem:sign-bridge}
and uses the same cyclotomic norms as the binary case.

\subsection{Sign correlation identities and exact bounds}

Put $\rho_s=\sum_i x_ix_{i+s}$. Expansion and orthogonality give
\[
 \rho_s=55-4k+4c_s,\quad \rho_s\equiv3\pmod4,\quad
 55+2\sum_{s=1}^{27}\rho_s=(2k-55)^2.
\]
For $p_r=|f_x(\zeta^r)|^2=4q_r$ and $\sigma=2k-55$,
\[
 \sigma^2+2\sum p_r=55^2,\qquad
 \sigma^4+2\sum p_r^2=55\left(55^2+2\sum_{s=1}^{27}\rho_s^2\right).
\]
At fixed weight this is an invertible affine change of the integer
correlations. The sign congruence follows from the binary
correlation variables, which are used in the exact profile enumeration.

The first-moment bound is
\[
 R(a)\le(55-2k)\left(\frac{k(55-k)}{54}\right)^{27}.
\]
The fourth-moment and capped-product lemmas above apply to $\prod q_r$,
with every threshold multiplied by $55-2k$. All exclusions are strict
integer or rational outward comparisons against the exact incumbent.
For a mod-$5$/mod-$11$ fold pair $(r,s)$, put
\[
 T_{55}=\frac{55k+k^2-5\sum r_i^2-11\sum s_i^2}{2}.
\]
The joint-fold bound is
\[
 (55-2k)N_5(r)N_{11}(s)(T_{55}/20)^{20}.
\]
The primitive fourth moment subtracts the two folded fourth moments
from the full fourth moment, so the same product lemma can be used
in that sector. The recorded joint screen combines the first-sector
bound with the global fourth-moment bound; all conservative survivors
are subsequently resolved by complete profiles and lifts.


\subsection{Complete profile and lift enumeration}

The exact moment screen leaves $k=22,23,24,25,26$. Capped moments
reduce $9,531,797,706$ ordered formal correlations to $731,378,844$
explicit evaluations in $17$ multisets.
\begin{center}
\small
\begin{tabular}{@{}rrrrr@{}}
\toprule
$k$ & Formal profiles & Explicit profiles & Targets & Lift tasks\\
\midrule
22 & 17,383,860 & 17,383,860 & 0 & 0\\
23 & 3,580,859,880 & 259,224,030 & 1,705 & 10,580\\
24 & 5,501,039,895 & 435,495,645 & 10,647 & 77,862\\
25 & 432,505,269 & 19,266,507 & 5,477 & 63,300\\
26 & 8,802 & 8,802 & 6 & 30\\
\bottomrule
\end{tabular}
\end{center}
There are $17,835$ canonical targets and $151,772$ lift tasks.
The weight-$22$ target set is empty.

The independent arithmetic checker parametrizes the two stationary
values by their separation, verifies its own dyadic cosine enclosures,
and generates multisets by descending multiplicities.
An independent bounded-alphabet FKM necklace generator produces
the cyclic fold representatives using a separate affine
canonicalization routine. The two fold lists and their
correlations agree exactly.

An independent \texttt{next\_permutation} traversal reconstructs
each canonical algebraic product by CRT and compares it exactly
with the incumbent. Its above-threshold set contains $255$
additional indefinite formal spectra. For every one, a recorded frequency and exact rational
upper bound prove $q_j<0$, excluding a binary autocorrelation.
After these explicit positive-semidefinite discharges, the independent
and production target sets, including their products, are exactly equal.

The production five-plus-six lift uses a sorted packed integer key.
The independent six-plus-five implementation uses typed seven-coordinate
hash keys with equality checks and tests every shift. Both exhaust all
$151,772$ tasks and inspect $113,997,335,580$ joined words.
Their joined counts and complete output sets agree task by task,
returning exactly two normalized words.

For global coverage, map any feasible word's correlation to its
canonical unit representative. Its two folds then occur in the
independent necklace lists in that orientation. The two fold rotations
are simultaneously induced by one translation modulo $55$, by CRT.
Consequently the Cartesian products of listed margin fibers contain
a transformed representative of every feasible word. An independent
task audit reconstructs all $151,772$ descriptors from those fold lists
and checks their exact partition. The two full lift enumerations
therefore prove global exhaustiveness.


\subsection{Sign extremizers and all equality cases}

The two normalized sign lift outputs are affine-equivalent to
the weight-$23$ representative $v$ defined in
Theorem~\ref{thm:sign}. Its sign encoding $x=2v-\mathbf1$ is
\begin{lstlisting}
----------+--++-++-+--+++--+--+++---+-+-+-++---+--+-+++
\end{lstlisting}
Its exact norms are
\[
 N_5=131,\quad N_{11}=917467,\quad N_{55}=15268987821561877723321,
\]
and its binary determinant is
\[
 \det C(a)=42208491650900637654045811643791.
\]
Direct integer Bareiss on the sign matrix gives a negative determinant
with absolute value
\[
 297532404423965431213849494795385191907933028352.
\]
Modular Gaussian elimination over three independently checked primes
and signed CRT agree, with sufficient capacity proved by Hadamard's
bound. Both agree with $2^{54}(2k-55)N_5N_{11}N_{55}$.

The affine orbit of $x$ has $2200$ distinct words and trivial
stabilizer. Its global negatives form a second orbit with $32$
plus signs and row sum $9$. One binary representative of this orbit is
\begin{lstlisting}
0000001101110110101010100110111101111011110011010001111
\end{lstlisting}
The two orbits contain $4400$ words in total.
Together with the exhaustive lift coverage, this proves
Theorem~\ref{thm:sign}.


\subsection{Comparison of the two maxima}

The sign representative $v$ and the binary maximizer $w$
occupy distinct affine classes.  The binary determinant
of $v$ falls below $D_{01}(55)$ by
$92485602443857757677261299685341$. Its weight-$32$ complement has
binary determinant $58724857949079148040411564026144$.
The half binary correlations are $9$ seventeen times and $10$ ten times;
the $54$ nontrivial sign correlations are $-1$ thirty-four times and
$3$ twenty times. Relative to the sign encoding of $w$, the nontrivial
Fourier product decreases, while the absolute sign row sum
increases from $1$ to $9$. The sign objective improves by approximately
$3.43$, the ratio of the two exact normalized integers.

At orders $15$ and $21$, direct Gray-code sign Fourier enumeration
covers all $2^{15}$ and $2^{21}$ words. The reduced sign moment/profile
pipeline followed by a CRT-column split lift agrees on the maximum,
weight reduction, profile coverage, lift coverage and complete winner
sets. The normalized maxima are $23859$ and $39337984$, with $480$
and $84$ winning words. A further $707$ exact bridge and moment checks
include all words at $3,5,7,9$ and one sampled word at every reduced
weight at order $55$.


\section{Conclusion}

The change from binary to sign circulants preserves every
nontrivial Fourier magnitude and changes the factor at frequency
zero. At order $55$, this changes the maximizing weight and
the affine-orbit classification. Cyclotomic folding isolates
the low-dimensional arithmetic constraints, while exact
correlation enumeration and binary lifts certify both extrema.
This provides a common proof structure for the two alphabet
variants.

\appendix
\section{Small-weight binary strata}\label{app:small}


Exact fixed-weight binary determinants for $1\leq k\leq7$
are listed below. Their complementary strata follow from
\eqref{eq:weight-pair}.
\begin{center}
\small
\begin{tabular}{@{}rr@{}}
\toprule
$k$ & $D_{\mathrm{circ}}(55,k)$\\
\midrule
1 & 1\\
2 & 2048\\
3 & 1564031349\\
4 & 6424318816256\\
5 & 12556202952818275\\
6 & 4342461217049330166\\
7 & 570999857161750276477\\
\bottomrule
\end{tabular}
\end{center}

\section{Exact certificate protocol and reproducibility}
\label{app:verification}

Both extremal determinant values are verified by three
exact computations: finite-field Fourier reconstruction using
signed CRT, integer fraction-free elimination, and the product
of the three cyclotomic resultants.  The profile and fold
screens use outward rational bounds for every pruning decision.

The binary verifier re-derives rational product bounds and cosine
enclosures, checks fold norms by resultants, regenerates
correlation profiles, and verifies lift tasks, output words,
determinants and affine orbits.  Its full-lift mode executes
the independent six-plus-five implementation.
Certificate manifests record input hashes, integer counts
and output identities.


The certificate package, source files and verifiers are at
\url{https://github.com/dongxuelian2/order55-binary-circulant-maxdet}.
From the repository root, the checks are:
\begin{lstlisting}
python scripts/verify_order55_global.py
python scripts/verify_order55_sign.py
python -m pytest
python -m pip check
\end{lstlisting}
Complete independent lift replays are requested with:
\begin{lstlisting}
python scripts/verify_order55_global.py --full-lifts
python scripts/verify_order55_sign.py --replay
\end{lstlisting}
Native replay requires a C++20 compiler with unsigned
\path|__int128| support. The binary companion report records
all $16,084$ lift tasks. The sign package records $17,835$
canonical targets, $151,772$ lift tasks, and $113,997,335,580$
joined words.

\bibliographystyle{unsrt}
\bibliography{references}

\end{document}
